\chapter{The forecast case}\label{ch:case}

A model like WRF is a large set of options; a forecast uses a small subset of them. The port, and most of this book,
is defined by one case: a coupled weather and fire simulation of the Eaton Fire, which started in the foothills above
Altadena, California, on the evening of 7~January~2025 local time. The case's files are in
\code{cases/eaton_20250108/}. The ignition used in the case is at 34.18604~N, 118.09325~W, at 02:18~UTC on
8~January~2025.

\section{Domains and time steps}

Two domains with one-way nesting (\code{feedback = 0}, \code{smooth_option = 0}). The ratio is 9 in both space and
time.

\begin{center}
\begin{tabular}{lrrrrr}
\toprule
Domain & Grid & $\Delta x$ & $\Delta t$ & Steps in 17 h & Fire grid \\
\midrule
d01 & $450\times450\times60$ & 900 m & 3 s & 20{,}400 & -- \\
d02 & $811\times811\times60$ & 100 m & 1/3 s & 183{,}600 & $3244\times3244$ (25 m) \\
\bottomrule
\end{tabular}
\end{center}

Each large time step has three Runge--Kutta stages, with 1, 2 and 4 acoustic sub-steps
(\code{time_step_sound = 0} gives four per step): seven acoustic sub-steps in all (\cref{ch:time}). Radiation is
called every minute (\code{radt = 1}), so every 20 steps of d01 and every 180 steps of d02.

\section{Inputs}

\code{real.exe} of WRF~4.6.0 made the three input files, with the MODIS land use of 21 categories
(\code{MODIFIED_IGBP_MODIS_NOAH}):
\begin{center}
\begin{tabular}{lrl}
\toprule
File & Size & Content \\
\midrule
\code{wrfinput_d01} & 289 MB & initial state of d01 \\
\code{wrfinput_d02} & 876 MB & initial state of d02, with the fire-grid fuel categories (LANDFIRE) and terrain \\
\code{wrfbdy_d01} & 122 MB & lateral boundary conditions of d01, every 3 h, 00:00--21:00 UTC \\
\bottomrule
\end{tabular}
\end{center}
Their md5 checksums, and those of the eight lookup tables the run reads, are in \code{cases/eaton_20250108/manifest.md5}:
\begin{itemize}
  \item the Noah tables \code{LANDUSE.TBL}, \code{VEGPARM.TBL}, \code{SOILPARM.TBL}, \code{GENPARM.TBL};
  \item \code{RRTMG_LW_DATA};
  \item the ozone climatology files \code{ozone.formatted}, \code{ozone_lat.formatted}, \code{ozone_plev.formatted};
  \item the greenhouse-gas table \code{CAMtr_volume_mixing_ratio}.
\end{itemize}
A copy of an input with a different checksum is a different input, and is never mixed into a comparison.

There is no \code{namelist.fire}. WRF-Fire therefore uses its built-in 53-category fuel table and writes it out as
\code{namelist.fire.output}.

\section{Options}

\paragraph{Dynamics.}
\begin{itemize}
  \item Non-hydrostatic, with the hybrid vertical coordinate (\code{hybrid_opt = 2}) and moist potential
    temperature (\code{use_theta_m = 1}); the hypsometric option \code{hypsometric_opt = 2}.
  \item Advection: fifth order in the horizontal, third order in the vertical; positive-definite for moisture and
    TKE (\code{moist_adv_opt = 1}, \code{tke_adv_opt = 1}).
  \item Turbulence:
    \begin{itemize}
      \item d01: two-dimensional Smagorinsky horizontal diffusion with the YSU boundary layer (\code{km_opt = 4});
      \item d02: a three-dimensional TKE large-eddy simulation (\code{km_opt = 2}), with sub-grid stress output
        (\code{m_opt = 1}).
    \end{itemize}
  \item Damping: Rayleigh damping of $w$ at the top (\code{damp_opt = 3}) and vertical-velocity damping
    (\code{w_damping = 1}); off-centering \code{epssm} of 0.5 on d01 and 0.8 on d02.
\end{itemize}

\paragraph{Physics.}
\begin{itemize}
  \item WSM6 microphysics (\cref{ch:mp}).
  \item Radiation (\cref{ch:rad}), called every minute:
    \begin{itemize}
      \item RRTMG longwave, in single precision in this build;
      \item Dudhia shortwave;
      \item the ozone climatology (\code{o3input = 2}) and greenhouse gases from a table (\code{ghg_input = 1}).
    \end{itemize}
  \item The revised MM5 surface layer (\cref{ch:sfc}) and the Noah land-surface model (\cref{ch:lsm}).
  \item YSU boundary layer on d01 only (\cref{ch:pbl}); none on d02, which resolves the boundary layer.
  \item No cumulus, shallow-convection, urban or lake schemes.
\end{itemize}

\paragraph{Fire, on d02.}
\begin{itemize}
  \item WRF-Fire (\code{ifire = 2}), refinement 4 in each direction.
  \item Hybrid WENO5/ENO1 level-set advection (\code{fire_upwinding = 9}), with re-initialization.
  \item Wind interpolated to 1~m with a logarithmic profile (\code{fire_lsm_zcoupling}).
  \item Heat and moisture fluxes fed back into the atmosphere (\code{fire_atm_feedback = 1}).
  \item The built-in fuel table and a uniform fuel moisture of 8\,\%.
  \item One circular ignition of radius 100~m, between 8280~s and 9080~s of model time.
\end{itemize}

\paragraph{Output.} History of d02 every 15~minutes (69 frames in 17~h); no history of d01; restarts every hour for
both domains.

\section{The reference runs}

The original run was made on the CCR cluster with the stock WRF build. The port compares against a different
reference, CPU-REF:
\begin{itemize}
  \item the same source built with flags that make the CPU arithmetic reproducible, and with a portable library of
    elementary functions (\cref{ch:repro});
  \item run on the host CPUs of the same workstation that holds the GPUs, in the same container image, so that the
    comparison is between two builds of one source on one machine; nothing is developed or tested on CCR
    (\code{roadmap/decisions/ADR-006-machines.md}).
\end{itemize}
CPU-REF's 17-hour run, its traces and its hourly restarts are archived with checksums. How far CPU-REF moves the
fire against the original run is measured once and recorded (\tbr); it is not a criterion, because no GPU can
reproduce another compiler's math library bit for bit.

\section{Development cases and windows}

A full d02 step costs minutes while most of the model still runs on one CPU core, so development uses smaller cases.
\begin{description}
  \item[S-3M] The WRF \code{em_fire} idealized case with the Eaton physics, three minutes of model time. It runs in
    minutes on one CPU core and needs no external data. It is the case of the CPU-only verification
    (\code{port/gates/cpu_verify.sh}).
  \item[eaton\_small] The real case with a $181\times181\times60$ nest centered on the ignition, and its own CPU-REF
    references.
\end{description}

Tests run in \emph{windows}: short runs that start from a restart file and compare a GPU run with a CPU-REF run that
started from the same file. The windows are defined in \code{port/h100/windows.txt}:
\begin{center}
\begin{tabular}{llrl}
\toprule
Window & Start & Length & Purpose \\
\midrule
W-T0 & 00:00 & 9 s & first d01 steps, nest opening \\
W-20 & 02:20 & 9 s & the everyday check of one routine \\
W-100 & 02:20 & 36 s & sub-phase gates \\
W-RAD & 02:20 & 240 s & four radiation calls per domain \\
W-TKE & 02:20 & 336 s & turbulence and sub-grid stresses \\
W-FORCE & 02:20 & 60 s & nest forcing \\
W-IGN & 02:00 & 35 min & the whole ignition phase \\
W-FIRE & 02:20 & 40 min & free spread \\
W-1H & 02:00 & 1 h & long-window check \\
\bottomrule
\end{tabular}
\end{center}

\section{What the case means for the port}

The options fix the code paths:
\begin{itemize}
  \item every routine the case executes in a time step must run on the GPU;
  \item every routine it does not execute may stay as it is.
\end{itemize}
A startup check in the GPU build (\code{gpu_check_config}) refuses a namelist outside the supported envelope. Dates,
sizes, the nest position and the ignition may change; the scheme choices may not. The full list of options, with the
reasons for each, is in \code{plan.md} section~2. \Cref{app:namelist} reproduces the namelist.

\paragraph{Memory.} The full case needs about 58--62~GB of device memory: the state of both domains, a scratch pool,
work arrays, the radiation batch, and the local memory of the column physics (\code{plan.md} section~3). It fits an
80~GB GPU, not one of the project's 40~GB A100s. The acceptance run of the port therefore uses \code{eaton_mid}, a
case with the same options and ignition and the largest nest that fits one 40~GB GPU with 15\,\% headroom (about
$400 \times 400 \times 60$, by the memory estimator); the full case follows on both A100s together
(\cref{ch:mgpu}). A smaller case still, \code{eaton_small} ($181 \times 181$ nest), serves development.
